%قبل از شروع مسئله ترجیحاً چیزی ننویسید. اما اگر خواستید بنویسید، در قالب لیستی (فقرات و شمارش) بنویسید.
\شروع‌مسئله[33] \متن‌سیاه{سؤال‌های جواب کوتاه.} در سؤال‌های این قسمت معمولاً یک جواب کوتاه و/یا یک توضیح کوتاه حدوداً بین یک تا دو خط مدنظر است؛ مگر این که صورت آن قسمت به صراحت خواستۀ دیگری داشته باشد.
% صورت اصلی مسئله را در این‌جا بنویسید.
\شروع{زیرمسائل}
\زیرمسئله[4]
در درهم‌سازی کامل، زمان یک جست‌وجو حداکثر چقدر می‌تواند باشد؟ چرا؟

\پاسخ{
$O(1)$.
هر جست‌وجو به دقیقاً دو بار محاسبۀ تابع درهم‌سازی و دو بار نگاه به خانه‌های جدول و مقایسه نیاز دارد.
	
}
\زیرمسئله[4]
اگر ضرب دو عدد را با الگوریتم فیور در زمان
$\Theta(d \log d\cdot2^{O(\log*d)})$
 پیاده‌سازی کنیم که $d$ تعداد رقم‌های عدد است، زمان تقسیم دو عدد با استفاده از روش نیوتون و الگوریتم فیور چقدر خواهد بود؟ چرا؟

\پاسخ{
	$\Theta(d \log d\cdot2^{O(\log*d)})$
	در کلاس دیدیم که با این روش، زمان تقسیم با ضرب به صورت مجانبی برابر می‌شود.
	
}
\زیرمسئله[4]
تبدیل باروز ویلر \متن‌تایپ{baba\$} چیست؟ نحوۀ به دست آوردن پاسخ را به اختصار توضیح دهید. 

\پاسخ{
\متن‌تایپ{abba\$}. کافی است چرخش‌های رشتۀ \متن‌تایپ{baba\$} را مرتب می‌کنیم و ستون آخر را درنظر بگیریم.
\begin{latin}
\texttt{\$bab\underline{\textbf{a}}\\}
\texttt{a\$ba\underline{\textbf{b}}\\}
\texttt{aba\$\underline{\textbf{b}}\\}
\texttt{ba\$b\underline{\textbf{a}}\\}
\texttt{baba\underline{\textbf{\$}}}
\end{latin}

}

\زیرمسئله[4]
%2013 Final P1.e
این جمله درست است یا غلط؟ چرا؟

روش وارسی درجۀ دو
(با استفاده از رابطۀ $h(x,i)=(h'(x) + a\cdot i^2+b\cdot i) \mod m$ با فرض ثابت بودن $a$ و $b$) 
در درهم‌سازی با آدرس‌دهی باز، شرط درهم‌سازی یکنواخت را برآورده می‌کند؛ یعنی دنبالۀ وارسی کلید $k$ به احتمال یکسان، برابر  با یکی از $m!$ جایگشت اعداد 0 تا $m-1$ خواهد بود.

\پاسخ{
غلط. تمام کلیدهایی که $h'(k)$ یکسان دارند، جایگشت یکسانی خواهند داشت. دلایل دیگر نیز درصورت درست بودن قابل قبول است.

}

\زیرمسئله[4]
به اختصار توضیح دهید با استفاده از آرایۀ پسوندی و آرایۀ کمکی طولانی‌ترین پیشوند مشترک (LCP) برای یک رشتۀ $s$، چطور می‌توان کوچکترین زیررشتۀ ناتهی $s$ را پیدا کرد که فقط یک بار در $s$ تکرار شده است؟

\پاسخ{
	تمام خانه‌های آرایۀ پسوندی را به صورت زیر بررسی می‌کنیم. اگر $SA$ آرایه پسوندی باشد و $SA[i]=j$ باشد، در این صورت کوچکترین رشتۀ یکتای $s$ که از $s[j]$ شروع می‌شود را می‌توان با کمک LCP و مقایسه با $s[SA[i-1]]$ و $s[SA[i+1]]$ در $O(1)$ پیدا کرد. در نهایت کمترین مقدار پیدا شده را خروجی می‌دهیم.
	
}

\زیرمسئله[4]
یک داده‌ساختار پشته داریم که علاوه بر \متن‌لاتین{pop()} (با هزینۀ $c$) و \متن‌لاتین{push(x)} (با هزینۀ $c$)، دارای عمل \متن‌لاتین{multipop(k)} (با هزینۀ $kc$) باشد که معادل $k$ بار اجرای \متن‌لاتین{pop()} است. می‌خواهیم با تحلیل تابع پتانسیل ثابت کنیم هزینۀ سرشکن هر عمل $O(1)$ است. یک تابع پتانسیل مناسب برای این کار پیشنهاد دهید.
چرا این تابع خوب کار می‌کند؟

\پاسخ{
	تعداد اعضای پشته ضرب در $c$. در صورتی که ضرب در $c$ لحاظ نشود، نیمی از نمره داده می‌شود.
	
}
\صفحه‌جدید
\زیرمسئله[4]
یک گراف وزن‌دار داریم که وزن یال‌ها می‌تواند یک عدد صحیح مثبت، $0$، $-1$ و یا $-2$ باشد. آیا می‌توان به وزن همۀ یال‌ها ۳ واحد اضافه کرد و با اجرای الگوریتم دکسترا کوتاهترین مسیرها از رأس مشخص $s$ به بقیه را پیدا کرد؟ اگر بله به اختصار توضیح دهید و اگر خیر یک مثال نقض ارائه دهید.

\پاسخ{
	خیر. گراف سه رأسی مسیر $P_3$ و رأس $t$ را به عنوان یک مقصد در نظر بگیرید. اگر برای وزن یال‌ها داشته باشیم $w(s,t)=-1$، $w(s,v)=0$ و $w(v,t)=-2$، در این صورت مسیر به درستی پیدا نمی‌شود.
	
}

\زیرمسئله[5]
یک درخت ریشه‌دار $T$ داریم و می‌خواهیم از روی آن داده‌ساختاری بسازیم که هر گاه دو رأس از $T$ به ما داده شد، بتوانیم پایین‌ترین جد مشترکشان را بیابیم. برای این کار روی درخت $T$، جست‌وجوی عمق اول (DFS) را اجرا می‌کنیم و هر موقع به یک رأس وارد می‌شویم (چه برای بار اول و چه وقتی از بچه‌های آن رأس برمی‌گردیم)، عمق رأس را در یک لیست یادداشت می‌کنیم. فرض کنید برای هر رأس اولین و آخرین اندیس مربوط به آن در لیست را داریم و برای هر خانۀ لیست هم می‌دانیم مربوط به کدام رأس است. با کمک این لیست چطور می‌توانیم به پرسش پایین‌ترین جد مشترک دو رأس پاسخ دهیم؟ از بین داده‌ساختارهایی که در این درس شنیده‌اید کدام را پیشنهاد می‌کنید؟ زمان ساخت داده ساختار و پاسخ به پرسش آن را بیان کنید.

\پاسخ{
	برای پاسخ به هر پرسش برای رأس $a$ و $b$ (با فرض این که اولین خانه مربوط به $a$ در لیست قبل از اولین خانه مربوط به $b$ در لیست است)، کافی است در زیربازۀ اولین اندیس مربوط به $a$ تا آخرین اندیس مربوط به $b$ کمترین عمق را پیدا کنیم که معادل «مسئلۀ پرسش کمینۀ یک محدوده» است که در جلسه ۱۷ مطرح شد و بهترین راه آن (که فقط نام برده شد)، در $O(n)$ ساخته می‌شود و پاسخ به هر پرسش در $O(1)$ انجام می‌شود. توجه کنید طول لیست با وجود این که هر رأس چند بار در آن ظاهر شده است با فرض $n>1$،  $2e-2 = 2n-4=O(n)$ هست و یعنی طول لیست با تعداد رأس‌ها از لحاظ مجانبی برابر است.
	
}

%\زیرمسئله[3]
%
%
%\پاسخ{
%
%
%	
%}
%
%\زیرمسئله[3]
%
%
%
%\پاسخ{
%
%
%	
%}

%%\begin{figure}[h!]%[5]{L}{0.4\textwidth}
%	%	\vspace{-2cm}
%	\begin{center}
%		\begin{tikzpicture}[-,>=stealth',level/.style={sibling distance = 5cm/#1,
%			level distance = 1cm}] 
%		\node [arn_s] {5}
%		child{ node [arn_s] {1} 
%			child{ node [arn_x] {}}
%			child{ node [arn_x] {}}
%		}
%		child{ node [arn_s] {15}
%			child{ node [arn_s] {10} 
%				child{ node [arn_s] {7}
%					child{ node [arn_x] {}}
%					child{ node [arn_x] {}}
%				}
%				child{ node [arn_s] {13}
%					child{ node [arn_x] {}}
%					child{ node [arn_x] {}}
%				}
%			}
%			child{ node [arn_s] {20}
%				child{ node [arn_x] {}}
%				child{ node [arn_x] {}}
%			}
%		}
%		; 
%		\end{tikzpicture}
%	\end{center}
%%	\caption{درخت قرمز-سیاه}
%%	\label{شکل:۷و۸}
%%\end{figure}



\پایان{زیرمسائل}



%اگر مسئله چند قسمت داشت به این صورت عمل کنید:
%\شروع{زیرمسائل}
%\زیرمسئله[0] % مسئلۀ ۱آ
%\زیرمسئله[0] % مسئلۀ ۱ب
%\زیرمسئله[0] % مسئلۀ ۱پ
%\زیرمسئله[0] % مسئلۀ ۱ت
%\پایان{زیرمسائل}





%%%%%%%%%%%%%%%%%%%%%%%%%%%%%%%%%%%%%%%%%%%%%%%%%%%%%
% برخی دستورات مفید برای تمرین‌ها. %
%%%%%%%%%%%%%%%%%%%%%%%%%%%%%%%%%%%%%%%%%%%%%%%%%%%%%

% فرمول‌های نمونه:
%$f(x) = \Theta(x)$
%$T(x, y) \leq \log(x) + 2^y + \binom{2n}{n}$
% {\tt code\_function}


% لیست موارد بدون شماره‌گذاری:
%\شروع{فقرات}
%    \فقره مورد اول 
%        \شروع{فقرات}
%            \فقره مورد اول داخلی 
%                \شروع{فقرات}
%                    \فقره اولین مورد خیلی داخلی 
%                    \فقره دومین مورد خیلی داخلی 
%                \پایان{فقرات}
%            \فقره مورد دوم داخلی 
%        \پایان{فقرات}
%    \فقره مورد دوم 
%\پایان{فقرات}

% لیست موارد شماره‌گذاری شده:
%\شروع{شمارش}
%    \فقره یک مورد 
%    \فقره یک مورد دیگر 
%    \فقره باز هم یک مورد دیگر
%\پایان{شمارش}

% برای نوشتن یک فرمول چند خطی که یک شماره بخوره. برای حالت‌های دیگر به سایت زیر مراجعه کنید:
% https://www.overleaf.com/learn/latex/Aligning_equations_with_amsmath
%\begin{align} 
%    \begin{split}
%        (x+y)^3 &= (x+y)^2(x+y) \\
%                &= (x^2+2xy+y^2)(x+y) \\
%                &= (x^3+2x^2y+xy^2) + (x^2y+2xy^2+y^3) \\
%                &= x^3+3x^2y+3xy^2+y^3
%    \end{split}                                 
%\end{align}

% نوشتن ماتریس و آرایه دوبعدی
%\begin{align}
%    A = 
%    \begin{bmatrix}
%        A_{11} & A_{21} \\
%        A_{21} & A_{22}
%    \end{bmatrix}
%\end{align}

% اضافه کردن تصویر با شرح. برای اضافه کردن تصویر بدون شرح، فقط دستور درج‌تصویر را استفاده کنید.
%\شروع{شکل}[h!]
%\تنظیم‌ازوسط
%  \درج‌تصویر[پهنا=0.5\پهنای‌متن]{img.jpg}
%\شرح{شرح تصویر}
%\برچسب{تصویر:۱}
%\پایان{شکل}
%\پایان{زیرمسائل}

%کد زیر را ببینید.
%%\lstinputlisting{source_filename.py}
%\begin{latin}
%\begin{lstlisting}
%# Import the modules
%import sys
%import random
%
%ans = True
%
%while ans:
%question = raw_input("Ask the magic 8 ball a question: (press enter to quit) ")
%
%answers = random.randint(1,8)
%
%if question == "":
%sys.exit()
%
%elif answers == 1:
%print "It is certain"
%
%elif answers == 2:
%print "Outlook good"
%
%elif answers == 3:
%print "You may rely on it"
%
%elif answers == 4:
%print "Ask again later"
%
%elif answers == 5:
%print "Concentrate and ask again"
%
%elif answers == 6:
%print "Reply hazy, try again"
%
%elif answers == 7:
%print "My reply is no"
%
%elif answers == 8:
%print "My sources say no"
%\end{lstlisting}
%\end{latin}
